\chapter{初抵：樟宜机场与第一印象}

\section{出发与到达}

今年一月，我有机会参加在新加坡举行的 AAAI 国际会议。这是一个已有数十年历史的人工智能大会。过去很长一段时间里，它在计算机领域并不算最耀眼的学术平台，但随着近年来人工智能的迅速崛起，这个会议的分量也与日俱增。

我的行程从 1 月 17 日晚间启程，自纽约肯尼迪机场出发，中途经由阿布扎比转机，再飞往新加坡。辗转十余小时之后，终于在 1 月 19 日清晨抵达目的地。

\section{樟宜机场：效率的第一课}

飞机降落在樟宜机场。令我印象深刻的，是入境手续之简便。我经由机场的 APEC 通道进入——APEC（亚太经济合作组织）为成员经济体的商务旅行者提供专属快速通关卡，持卡人无需排队经过普通入境柜台——仅需将护照在读卡机上扫描一下，整个过程大约只有十秒钟，随即便顺利通关。与许多国家繁琐冗长的入境程序相比，这种效率令人颇感意外。航班于上午九点半落地，不到十点，我们便已走出航站楼。

这也许是新加坡给我上的第一堂课——效率，不是口号，而是嵌入日常运转之中的一种制度习惯。

\section{入城：Uncle 与那座城市的轮廓}

走出航站楼之后，我们用手机召了一辆 Grab。这是新加坡最为普及的网约车服务，不多时便有司机接单——一位华人司机。

在新加坡，我初次接触到一种颇有意思的称呼方式：若司机是华人男性，人们通常称其为"Uncle"；若是女性，则称"Auntie"。这种称呼带着一种家庭式的亲近感，令人莞尔。不过后来我才了解到，这一习惯主要流通于华人社群之间，对于马来族或印度族司机，一般并不如此相称。

那天正值周日早晨，路上出奇地安静。樟宜机场位于新加坡岛的东端，我们的车沿着滨海高速公路一路向城市中心驶去，几乎畅行无阻。沿途，司机 Uncle 兴致勃勃地为我们介绍窗外的景观：那片新兴的商业区，那座拔地而起的金融中心，还有一座座规模庞大、统称为"Mall"的购物综合体。

约莫二十五分钟之后，我们进入了新加坡沿海的核心城区。

\section{地标：一座城市的自我表达}

首先映入眼帘的，是一座巨大的圆形建筑——新加坡摩天观景轮。它的轮廓在晨光中格外醒目，如同一枚悬于海边的巨型钟表。

稍远处，三座高楼并肩而立，顶端横跨着一条形如船舶的弧形结构。这便是举世闻名的滨海湾金沙——新加坡全国仅有的两家合法赌场之一便设于其中。它的存在，本身便是一个耐人寻味的政策选择：这个以严格法纪著称的国家，却以两处赌场作为其城市天际线的核心地标。

再远处，是一片开阔的滨海公园，绵延至海。

\section{Andaz 酒店：城市的垂直逻辑}

车子离开海滨，驶入城市腹地，在几条蜿蜒的街道间穿行片刻，停在了一座高楼入口处。这便是我们此行入住的 Andaz 酒店。

这座酒店并非一栋独立建筑，而是嵌于一个大型综合开发项目之中。进入建筑后，电梯直接从入口层升至二十五层——那里才是酒店的大堂。换言之，二十五层以下，是各类商业空间、办公区域及其他用途；二十五层以上，才属于酒店的领地。

这种垂直叠加的空间利用方式，在寸土寸金的新加坡并不罕见。它折射出这座城市面对土地约束时的一贯应对逻辑：向上生长。

\section{一个悬而未决的问题}

我们在新加坡共停留了十天。大部分时间自然用于参加会议和处理工作事务，但因酒店距市中心不远，每日清晨与傍晚，我仍能抽出一些零散时间，在城中漫步。

在这些独处的间隙里，我反复思索一个问题：

\begin{center}
\textit{为什么在中国改革开放初期，一个幅员辽阔的大国，会选择向新加坡这样一个弹丸小国取经？}
\end{center}

这让我想起一段流传颇广的往事。据说，在李光耀第一次访问北京时，邓小平曾问他：中国能不能成为一个发达国家？

李光耀后来回忆说，当年留守北京的是中国的精英阶层，而漂洋过海来到新加坡的华人，大多是社会底层的劳工。如果这些出身贫苦的移民，能够在一座几乎毫无自然资源的小岛上建立起现代化的国家，那么中国当然也有理由实现现代化。

在这个意义上，新加坡不仅仅是一座城市国家，它更像是一场持续进行中的实验——一个小国，如何在强权林立的世界体系中，找到属于自己的生存之道。
